\section{Boolean ancestry, coverage, and collision}\label{sec:duality}

This section is valid for every positive loopless row-stochastic $P$;
reversibility is not used.

\subsection{The fair-geometric union dual}

At fitness two, \eqref{2.2} becomes $2x/(1+x)$.  This is a finite additive
graphical system in the sense of \citet{Harris1978}.  Let $L$ have the fair
geometric law
\[
 \Pp(L=\ell)=2^{-\ell},\qquad \ell\ge1,                       \tag{3.1}\label{3.1}
\]
and, conditionally on target $v$, draw $U_1,\ldots,U_L$
independently from row $P_v$.  For any mutant set $S$,
\[
 \Pp(U_1,\ldots,U_L\notin S)
 =\E(1-P_{vS})^L=\frac{1-P_{vS}}{1+P_{vS}}.                  \tag{3.2}\label{3.2}
\]
Thus the forward update at $v$ is exactly the additive Boolean map
\[
 x_v\longleftarrow x_{U_1}\vee\cdots\vee x_{U_L}.           \tag{3.3}\label{3.3}
\]
The transpose graphical map sends a nonempty ancestral set $A$ to
\[
 A\longmapsto
 \begin{cases}
 (A\setminus\{v\})\cup\{U_1,\ldots,U_L\},&v\in A,\\
 A,&v\notin A.
 \end{cases}                                                 \tag{3.4}\label{3.4}
\]
Because $P_{vv}=0$, the updated set in the first line omits $v$ and is
always nonempty; hence the dual remains a proper nonempty subset whenever it
starts there.

Using the same graphical marks forward and backward gives the exact duality
\[
 \Pp_S(X_t\cap A\ne\varnothing)
 =\Pp_A(S\cap A_t\ne\varnothing).                            \tag{3.5}\label{3.5}
\]
For positive $P$, the dual is irreducible on the proper nonempty subsets.
Let $\Pi_P$ denote its stationary law and put
\[
 m(P)=\E_{\Pi_P}|A|.                                         \tag{3.6}\label{3.6}
\]

\begin{proposition}[coverage representation]\label{prop:coverage}
The fixation committor at fitness two is
\[
 h_P(S)=\Pp_{A\sim\Pi_P}(A\cap S\ne\varnothing),            \tag{3.7}\label{3.7}
\]
and therefore
\[
 \rho_{\dB}(P,2)=\frac{m(P)}n.                              \tag{3.8}\label{3.8}
\]
Moreover, for every nonempty $T$ disjoint from $S$,
\[
 (-1)^{|T|+1}\Delta_T h_P(S)
 =\Pp_{\Pi_P}(A\cap S=\varnothing,\ T\subseteq A)\ge0.      \tag{3.9}\label{3.9}
\]
Thus $h_P$ is a normalized completely alternating coverage function.
\end{proposition}

\begin{proof}
The finite forward chain absorbs almost surely.  Start the dual from $V$ in
\eqref{3.5} and let $t\to\infty$.  The left side converges to the fixation
probability from $S$, while the recurrent dual law converges to $\Pi_P$,
giving \eqref{3.7}.  Averaging $h_P(\{i\})=\Pp(i\in A)$ over $i$ proves
\eqref{3.8}.  For fixed $A$, the indicator
$\one_{\{A\cap S\ne\varnothing\}}$ has mixed difference
$(-1)^{|T|+1}\Delta_T$ equal to
$\one_{\{A\cap S=\varnothing,\ T\subseteq A\}}$; averaging proves
\eqref{3.9}.
\end{proof}

\subsection{A one-sample active chain}

The geometric burst can be factored into one-sample steps.  Let
\[
 \mathcal Y_n=\{(B,v):\varnothing\ne B\subseteq V\setminus\{v\}\}. \tag{3.10}\label{3.10}
\]
From $y=(B,v)$, with $b=|B|$, define $K(P)$ by the following fair mixture:
\begin{enumerate}
\item sample $i$ from row $P_v$ and move to $(B\cup\{i\},v)$;
\item choose $w$ uniformly from $B$, sample $i$ from row $P_w$, and move to
      $((B\setminus\{w\})\cup\{i\},w)$.
\end{enumerate}
Each branch has probability $1/2$.  The first branch continues a geometric
burst; the second stops it and retargets the active ancestral mark.

\begin{lemma}[active collision identity]\label{lem:active}
Let $\nu(P)$ be the stationary law of $K(P)$ and put
\[
 H(B,v)=\frac1{|B|}.                                        \tag{3.11}\label{3.11}
\]
Then
\[
 \nu(P)H=\frac1{m(P)}=\frac1{n\rho_{\dB}(P,2)}.             \tag{3.12}\label{3.12}
\]
The map $P\mapsto K(P)$ is linear.
\end{lemma}

\begin{proof}
We include the marked-current calculation because it is the bridge from the
coverage law to the Hessian.  Let $T_v(A,B)$ be the transition kernel of one
union-dual update with target fixed at $v$.  For $v\notin C$, define
\[
 \sigma_v(C)=\Pi_P(C\cup\{v\}),\qquad
 \eta_v(C)=\sum_A\Pi_P(A)T_v(A,C)-\Pi_P(C).                 \tag{3.13a}\label{3.13a}
\]
The subtracted term is exactly the passive current from $A=C$, for which
$v\notin A$ and no ancestral update occurs; hence $\eta_v(C)$ is the
nonnegative active current arriving at output $(C,v)$.  Stationarity of the
union dual, after summing its $n$ equally likely target updates, gives
\[
 \sum_{v\notin B}\eta_v(B)=|B|\Pi_P(B).                    \tag{3.13}\label{3.13}
\]
If $A_v$ adjoins one sample from row $P_v$, the fair-geometric resolvent
identity gives
\[
 \lambda_v=\frac{\sigma_v+\eta_v}{2},\qquad
 \eta_v=\lambda_vA_v.                                      \tag{3.14}\label{3.14}
\]
Write $A$ for the one-sample channel and $R$ for the fair
continue-or-retarget channel.  After the sample, continuing contributes
$\eta_v(B)/2$ to $(B,v)$.  Stopping, choosing $w$ uniformly from $B$, and
ending at $(B\setminus\{w\},w)$ contributes, by \eqref{3.13},
$\Pi_P(B)/2=\sigma_w(B\setminus\{w\})/2$.  Hence
$\lambda AR=\lambda$.  Since $\eta=\lambda A$, associativity gives
$\eta RA=\eta$; the push-forward $RA$ is precisely the active kernel
$K(P)$ defined above.  Both $\sigma_v$ and $\eta_v$ have total mass
$\Pp_{\Pi_P}(v\in A)$, so
\[
 \sum_{v,C}\lambda_v(C)=m(P).                              \tag{3.15}\label{3.15}
\]
Normalizing the pushed-forward law gives $\nu(P)=\eta/m(P)$.  Dividing
\eqref{3.13} by $|B|$ and summing over $B$ yields
\[
 \nu(P)H=\sum_{k\ge1}\frac{k\Pi_P(|A|=k)}{m(P)}\frac1k
 =\frac1{m(P)}.                                             \tag{3.16}\label{3.16}
\]
Both active moves contain exactly one use of a row of $P$, proving linearity.
\end{proof}

At the complete kernel, put $N=n-1$.  Direct substitution gives
\[
 \nu_0(B,v)=\frac{|B|}{nN2^{N-1}},\qquad
 c_0:=\nu_0H=\frac{2^N-1}{N2^{N-1}}=\frac1{m_n}.             \tag{3.17}\label{3.17}
\]
This agrees with \eqref{2.4} at $r=2$.

The collision interpretation is literal.  In one stationary marked-chain
step, let $I$ be the sampled source; if the fair continue-or-stop coin stops,
let $W$ be the new target chosen uniformly from $B$.  Then
\[
 \Pp(\mathrm{stop},W=I)=\frac1{2m(P)},\qquad
 \Pp(W=I\mid\mathrm{stop})=\frac1{m(P)}.                  \tag{3.18}\label{3.18}
\]
Thus maximizing fixation is equivalent to minimizing this stationary
collision rate.
